\documentclass[ecta,nameyear,draft]{econsocart}
\usepackage{amsmath,amsfonts,amssymb,amsthm,booktabs,array,longtable}
\usepackage[T1]{fontenc}
\usepackage{mathpazo,microtype,needspace,xurl}
\usepackage{xr-hyper}
\RequirePackage[colorlinks,citecolor=blue,linkcolor=blue,urlcolor=blue]{hyperref}
\startlocaldefs
\newtheorem{theorem}{Theorem}
\newtheorem{proposition}[theorem]{Proposition}
\newtheorem{lemma}[theorem]{Lemma}
\newtheorem{corollary}[theorem]{Corollary}
\newcommand{\E}{\mathbb E}
\newcommand{\R}{\mathbb R}
\endlocaldefs
\setlength{\emergencystretch}{3em}
\makeatletter\def\copyright@text{Prepared for referee review}\makeatother

\setcounter{secnumdepth}{2}

\externaldocument[][nocite]{build/supp}[supp.pdf]
\externaldocument[][nocite]{build/appendix}[appendix.pdf]
\begin{document}
\begin{frontmatter}
\title{Neural Bellman Operators}
\runtitle{Neural Bellman Operators}
\begin{aug}
\author[id=au1,addressref={add1}]{\fnms{Qian}~\snm{QI}\ead[label=e77ECTA]{qiqian@pku.edu.cn}}
\address[id=add1]{Peking University}
\end{aug}
\begin{abstract}
This paper develops Neural Bellman Operators for constructive policy evaluation and feasible improvement. Attained continuations retain action witnesses and an original-optimum error account. We derive representation-calibrated residual qualification and exact action-switch certificates that transfer policy-cost conclusions to classes of optimization outputs. A prospective nonlinear-investment study compares nine complete controllers under the original continuous law. Calibrated qualification improves both ReLU and Bernstein constructions, while reversing their earlier matched cost ordering. Independently implemented whole-bin evaluation identifies the gains and conditional purchase regions. Three fixed-label optimization problems have complete action-stable minimizer classes, which inherit their attained-policy intervals without further sampling. Recursive query, vector-control, predictable evaluation and deployment-allocation results retain their assumptions and evidence. The numerical comparison identifies the effect of specified constructions and qualification rules, not universal neural dominance.
\end{abstract}
\begin{keyword}
\kwd{Dynamic programming}\kwd{Neural approximation}\kwd{Policy certification}
\kwd{Information acquisition}\kwd{Computational economics}
\end{keyword}
\end{frontmatter}
\section{Introduction}\label{sec:introduction65}
A continuation approximation matters because it produces an economic decision. Prediction error, original-optimum policy loss and the resources required to verify that loss are different objects. This paper connects them through a constructive Neural Bellman Operator (NBO), its attained own-future continuation and the complete feasible controller that is implemented under the original economic law.

The operator first fits each date against its own attained future and retains feasible action witnesses. Its affine--ReLU realization and primitive error budget do not require an unrestricted nonconvex optimizer to succeed. A complementary recursive implementation queries the original Bellman problem at the states needed by action search. Primitive curvature, verified endpoint bounds and protected conventional refinement determine its validity; a trained proposal changes only optional query locations. This distinction separates the constructive neural continuation from the performance of a neural warm start.

A qualified shallow continuation can be evaluated without the original-value descendant tree. Its residual oscillation and implemented greedy gaps control original-optimum policy loss, while exact conditional feature integration supplies the evaluation primitives. Fixed-label qualification readouts have finite optimization certificates. Purchased Bellman information can also be reused subject to verified state displacement, capacity repair, retention and precision contracts. These constructions are retained together with their actual work accounts; the absence of a tensor state table does not erase horizon dependence.

We address two additional questions about the attained continuation. First, does a qualification bound charge a common represented function consistently across bases? Coefficientwise bounds can differ even on identical affine functions. We construct derivative-control qualification that is exact on the affine intersection of the ReLU and Bernstein classes and retains the original residual and policy-loss account. The construction is convex and admits exact primal--dual checks. It does not assert that the nonlinear function classes or their numerical conditioning are identical.

Second, can the economic conclusion survive changes in the numerical optimizer's output? For each robust alternative action, a rational dual certificate bounds the qualification objective required to change the implemented choice. A strict lower barrier certifies the action throughout an entire objective sublevel. With a fixed terminal actor in the two-date construction, this yields a complete-policy and original-law cost statement. Partially certified covers instead give an explicit cost enlargement through their initial-law probability. This is an optimization-to-policy transfer, not a fitting-seed success probability.

The prospective factorial experiment fixes both bases, both qualification rules and both observation resolutions before evaluating nine complete controllers. Both own-future fits are constructed under the declared rule. Two independent whole-bin implementations and exact rational checks qualify the continuous-law cost observations. The calibration improves the attained ReLU and Bernstein policies, but their matched ranking reverses. The earlier ReLU advantage under coefficientwise qualification and the new Bernstein advantage under calibrated qualification are both retained. The comparison identifies an operational qualification mechanism rather than attributing a ranking to a representation name.

The attained-cost event also supports conditional economic choices. Joint marginal and paired intervals have sharp difference-constraint and deployment-allocation interpretations. Where action stability holds throughout a readout class, the same cost intervals and conditional charge regions apply to that class without a new union bound over optimizers. Monetary decisions still require a loss scale, deployment count and resource charges; the experiment does not invent those inputs or count simulation rows as economic deployments.

\paragraph{Relation to existing methods.}
Lipschitz extensions and distance transforms underlie the constructive envelope \citep{mcshane1934,felzenszwalb2012}. Approximate and safe policy iteration have established error analyses \citep{munos2003,pirotta2013,scherrer2015}, as does convex dynamic-programming approximation \citep{yang2020convex}. Learned warm starts and optimization-based control supply related numerical methods \citep{sakaue2023warm,sambharya2023,agrawal2020control}. We retain these connections without claiming that a common verifier establishes the incremental value of its neural proposal.

Value-function control variates and telescoping identities are established evaluation tools \citep{munos2006variance,jiang2016}. Time-uniform concentration and off-policy confidence sequences have their own theory \citep{howard2020time,karampatziakis2021}. Our original-law costs include the numerical qualification required to use such bounds for an implemented controller. Minimax-regret decisions have an established economic role \citep{manski2004}; shortest-path difference constraints are classical \citep{floyd1962}. The new transfer identifies when the paid attained-policy evidence remains valid for an entire qualification sublevel, using exact robust-action switch barriers rather than an unobserved distance to the optimizer.

\paragraph{Reading path.}
The article develops the original operator, attained-continuation qualification, optimization-stable policies, predictable evaluation and actual economic costs. The Technical Supplement retains the core proofs. The Numerical and Development Appendix compiles the recursive query, pre-query readout, vector-control and historical evidence sections without alteration, together with all new dual proofs and factorial records. Historical companion editions retain diffusions, recursive and endogenous preferences, temporal selves and games under their own hypotheses. Their preservation is not a claim that the present expected-cost identities apply without modification to those different objectives.

\section{The economic operator and its policy account}\label{sec:core49}
Consider dates $t=0,\ldots,T$ and an invariant compact state domain $X\subset\R^d$. At date $t$, the action belongs to a nonempty compact set $A_t(x)$. The transition is $F_t(x,a,Z_{t+1})$, the current cost is $c_t(x,a)$, and the terminal cost is $g(x)$. Write $B_0=1$ and $B_t=\prod_{j<t}\beta_j$, with $\beta_j>0$. An implemented policy includes its observation rule, action selection and feasible numerical realization, not just its stored parameters.

For a continuation $v$, define
\begin{equation}\label{eq:operator49}
 Q_t(v;x,a)=c_t(x,a)+\beta_t\rho_t(v(F_t(x,a,Z))),\qquad
 \mathcal T_tv(x)=\inf_{a\in A_t(x)}Q_t(v;x,a).
\end{equation}
Throughout the deterministic comparison, the evaluations $\rho_t$ are monotone and cash invariant on a common domain containing all functions and constant translates used below. They are therefore nonexpansive in the supremum norm. Continuous objectives, effective compact feasible correspondences and measurable minimizers, or measurable arbitrarily near-minimizers with their error charged, are assumed. The stipulated conditional recursion gives Bellman comparison against all admissible adapted controls. In the experiments $\rho_t$ is conditional expectation, costs are bounded and innovations have their original continuous uniform law. The direct-cost inference in Section~\ref{sec:decisions49} is specific to expectation; it is not automatically an inference theorem for nonlinear recursive utility or games.

A Neural Bellman Operator is a construction and evaluation procedure for \eqref{eq:operator49}. It forms targets using its \emph{own} fitted $f_{t+1}$, returns a neural continuation and a feasible policy, and records an economic error allowance. It does not receive the optimal continuation as a training label. Arbitrary trained networks can be evaluated by the first result below. The subsequent constructive backend specifies how a network and its actor are obtained without assuming successful nonconvex optimization.

\subsection{Action-dependent errors and complete-policy loss}
Let $H(a)$ be a true current cost and $\widehat H(a)$ its continuation approximation. If $\widehat a$ is $\delta$-optimal for $\widehat H$ on the same feasible action set, then
\begin{equation}\label{eq:center49}
 H(\widehat a)-\inf_aH(a)\le\delta+
 \operatorname{osc}_{a}\{\widehat H(a)-H(a)\}.
\end{equation}
Indeed, compare $\widehat a$ with an arbitrarily near-optimal true action and cancel the approximate objectives. The infimum limit gives the result. Additive continuation levels cancel; action-dependent discrepancies do not. Under a change in future policies, the continuation must be transported to the new economic problem before this inequality licenses reuse. The full development edition retains the original quantitative transport, recursive-domain and economic-application results. No transport bound is presumed valid after an unaccounted change of constraints or preferences.

\begin{proposition}[One-sided policy sandwich]\label{prop:sandwich49}
Suppose, uniformly in state,
\[
 f_t-\mathcal T_tf_{t+1}\le u_t,\qquad
 \mathcal T_t^\pi f_{t+1}-f_t\le v_t,
 \qquad \ell_T\le f_T-g\le u_T.
\]
Then the actual feasible policy satisfies
\begin{equation}\label{eq:sandwich49}
 0\le J_0^\pi-V_0^*\le
 \sum_{t<T}B_t(u_t+v_t)+B_T(u_T-\ell_T).
\end{equation}
The bound is invariant to date-specific constant shifts of the fitted continuations, provided the associated residual accounts are shifted consistently.
\end{proposition}
\noindent The proof is in Supplement, Section~\ref{supp:proof-core49-1}.


A sharper valid residual account for the same continuation and the same deployed actor sharpens knowledge of the policy. It does not retrain that policy or change its realized cost. Conversely, subtracting two upper bounds in \eqref{eq:sandwich49} does not identify the difference in actual costs. These distinctions govern both the historical comparisons and the new experiments.

\section{Constructing and compiling a feasible neural policy}\label{sec:construction49}
Assume the primitive objective has, for feasible state--action pairs, state modulus $A_t^Q=C_{x,t}+\beta_tF_{x,t}L_{t+1}$ and action modulus $D_t=C_{a,t}+\beta_tF_{a,t}L_{t+1}$. Let
\[
 d_H(A_t(x),A_t(z))\le K_t^A\|x-z\|_1,
 \qquad L_t=A_t^Q+D_tK_t^A,\qquad L_T=L_g.
\]
The Hausdorff condition and the primitive moduli imply that $\mathcal T_tf_{t+1}$ is $L_t$-Lipschitz whenever $f_{t+1}$ is $L_{t+1}$-Lipschitz. All moduli and numerical integration procedures are supplied by the economic primitives, not fitted from successful outcomes.

At state nodes $x_{t,i}$ with covering radius $h_t$, form effective feasible $k_t$-nets. Evaluate their own-future objectives with certified error at most $e_t$. Minimize the computed labels and retain the corresponding \emph{original} action $a_{t,i}$. A Borel repair $\Gamma_{t,i}(x)$ is feasible at $x$ and satisfies
\begin{equation}\label{eq:repair49}
 \|\Gamma_{t,i}(x)-a_{t,i}\|_1\le K_t^A\|x-x_{t,i}\|_1+\zeta_t.
\end{equation}
The allowance $\zeta_t$ concerns displacement; feasibility itself may not fail. Define
\begin{equation}\label{eq:cone49}
 f_t(x)=\min_i\{y_{t,i}+L_t\|x-x_{t,i}\|_1\},\quad
 i_t(x)=\min\arg\min_i\{y_{t,i}+L_t\|x-x_{t,i}\|_1\}.
\end{equation}
Return $\pi_t(x)=\Gamma_{t,i_t(x)}(x)$. Terminal labels approximate $g$ with error $e_T$ and use the same cone construction with modulus $L_g$.

\begin{theorem}[Own-future construction with feasible witnesses]\label{thm:construct49}
The functions in \eqref{eq:cone49} have exact affine--ReLU realizations and Lipschitz moduli $L_t$ independent of label error. The returned policy is feasible and measurable. It satisfies
\begin{align}\label{eq:enclosures49}
 -e_t\le f_t-\mathcal T_tf_{t+1}&\le D_tk_t+e_t+2L_th_t,\\
 \mathcal T_t^\pi f_{t+1}-f_t&\le e_t+D_t\zeta_t,\notag
\end{align}
and hence
\begin{equation}\label{eq:constructionloss49}
 J_0^\pi-V_0^*\le
 \sum_{t<T}B_t(D_tk_t+2e_t+2L_th_t+D_t\zeta_t)
 +B_T(2e_T+2L_gh_T).
\end{equation}
With effective covers, repairs and certified queries at the prescribed accuracy, allocating each nonnegative term a share of a positive economic tolerance gives a terminating finite construction. A fixed numerical precision or a capped finite experimental ladder need not contain that allocation.
\end{theorem}
\noindent The proof is in Supplement, Section~\ref{supp:proof-core49-2}.


The Lipschitz-extension principle is classical \citep{mcshane1934}. A native min-plus envelope and its exact ReLU circuit, using the same labels, witnesses and ties, are \emph{one method with two encodings}. They return the same function and policy. Neither Theorem~\ref{thm:construct49} nor a faster representation proves a uniquely neural capability or a training-success probability for Adam or L-BFGS. The mathematical increment is the complete feasible, own-future construction and implementation account; the empirical question is its work and policy value relative to conventional generators.

\subsection{Original-owner compilation on a tensor cover}
On a tensor cover $\mathcal X=\prod_{j=1}^d\{x_{j,0},\ldots,x_{j,n_j}\}$, put $S=\prod_j(n_j+1)$. For arbitrary labels $y_i$, including inconsistent inexact labels, compute at each grid node $v$
\[
 (z_v,o_v)=\min_{\rm lex,i}(y_i+L\|v-x_i\|_1,i).
\]
The index $o_v$ is the original label owner, not the transformed node $v$.

\begin{theorem}[Continuous original-owner identity]\label{thm:compile49}
For a closed tensor cell $C$ with corners $\mathcal V(C)$, and $x\in C$,
\begin{align}\label{eq:compile49}
 f(x)&=\min_{v\in\mathcal V(C)}(z_v+L\|x-v\|_1),\\
 (f(x),i_*(x))&=\min_{\rm lex,v\in\mathcal V(C)}
 (y_{o_v}+L\|x-x_{o_v}\|_1,o_v).\notag
\end{align}
Coordinatewise forward and backward relaxations compute $(z,o)$ in $O(dS)$ additions and comparisons and $O(S)$ storage. After locating the cell, a query uses $O(d2^d)$ arithmetic operations. Exact rational comparisons preserve all ties, including $L=0$. Using the original owner's repaired action preserves the uncompiled policy, its cost and its certificate exactly.
\end{theorem}
\noindent The proof is in Supplement, Section~\ref{supp:proof-core49-3}.


The nodal algorithm is the classical distance transform \citep{felzenszwalb2012}; preserving the original action through continuous off-grid queries is the policy obligation. If there are $J_{t,i}$ node actions and $M_t$ integration representatives, write $Q_t=M_t\sum_iJ_{t,i}$. Literal dense cone evaluation costs $O(\sum_t dS_{t+1}Q_t)$ in representation work. Compilation replaces this by $O(\sum_t[dS_{t+1}+Q_td2^d])$, plus cell location, primitive evaluation, integration, feasible repair, certificate and output. Rational bit lengths and state boxes crossing many cells incur additional costs. There is no dimension-free complexity assertion.

\subsection{Acquired states and conventional actors}
A $b$-bit coordinate sensor reports a dyadic cell and its midpoint $\widehat x$. On the unit cube, $\|x-\widehat x\|_1\le r(b)=d2^{-b-1}$. For the investment capacity $\kappa(x)$ below, clip the original action to the smallest capacity throughout that cell and round down to an action quantum $\Delta_t$. This is feasible at every state in the cell. Relative to an exact true-state repair, its displacement is at most $2K_t^Ar(b)+\Delta_t$. Transporting the witness cone from $\widehat x$ to $x$ adds at most $2L_tr(b)$, plus any certified selector objective error $\nu_t$. Thus the implemented-policy account adds
\begin{equation}\label{eq:acquire49}
 \sum_{t<T}B_t\{(2L_t+2D_tK_t^A)r(b)+D_t\Delta_t+\nu_t\}.
\end{equation}
No continuity of the discrete owner selector is assumed. At uncertain numerical comparisons the simulator retains all possible actors, rather than choosing a favorable one.

For conventional FVI the continuation is multilinear on its realized tensor grid. Its actual modulus is bounded by the largest coordinate edge slope. A nearest-node actor uses a feasible node minimizer and true-state repair. Let $h_t$ be the largest half-cell $\ell^1$ radius, $L_t^{\rm gr}$ the Bellman graph modulus, and $e_t$ the node query allowance. Its optimal residual upper bound is $D_tk_t+e_t+L_t^{\rm gr}h_t$. Its selected-policy upper bound is $e_t+E_t$, where
\[
 E_t=\sup_x\{y_{i(x)}+L_t^{\rm gr}\|x-x_{i(x)}\|_1-f_t(x)\}.
\]
Partitioning cells by nearest-node boundaries makes the expression separately affine in each coordinate. Its maximum is therefore at a partition vertex; the implementation encloses all such vertices. Both uniform and curvature-refined FVI receive this sharper one-sided account, not an artificially doubled baseline. Their acquired-state allowance is derived separately in \eqref{eq:fvisensor49}.

\section{Original-optimum accuracy in the investment economy}\label{sec:accuracy62}
The exact witness theorem answers an implementation question conditional on a continuation. An economist also needs to know how close the complete implemented policy is to the original optimum. We now establish that connection in the investment economy itself. The argument exploits a property not used by the earlier cell-range calculation: despite the bilinear transition, the Bellman objective remains jointly convex. This permits a second-order interpolation account and transfers verified nodal decisions to an implementable policy over the entire state domain. It does not assume that a fitted neural continuation is convex.

\paragraph{Relation to the preceding two-control query.}
The earlier query experiment used the fixed shared capacity $1/8$. The complete dynamic extension below uses the original state-dependent capacity $c(x)$, with the retained exposure columns, action-cost interaction and second shock. Its nesting claim concerns the original scalar economic model; it is not an assertion that a full dynamic policy was already present in the earlier isolated query.

\subsection{Primitives and Bellman regularity}
Let $d\ge2$ be even, let indices be cyclic, and set $\beta=15/16$. Write
\begin{align*}
 S(x)&=\frac2d\sum_j(x_j-5/8)^2+
       \frac1{4d}\sum_j(x_j-x_{j+1})^2+
       2\left(\frac12-\frac2d\sum_jx_j\right)_+^2,\\
 g(x)&=S(x)+\frac2d\sum_j(x_j-5/8)^2,\qquad
 c(x)=\frac18+\frac1{8d}\sum_jx_j.
\end{align*}
The original scalar control has $0\le a\le c(x)$, running cost $S(x)+a^2+4a^4$, and transition
\[
 F_j(x,a,z)=\frac1{16}+\frac{x_j}{2}+\frac{x_{j+1}}8+
       \frac{x_j(1-x_{j+1})}{16}+B_ja+s_jz,
\]
where $(B_j,s_j)=(1/2,1)$ at alternating coordinates and $(1/4,-1)$ at the others, and $z$ is uniform on $[-1/32,1/32]$.

For the retained two-control extension, $a_1,a_2\ge0$, $a_1+a_2\le c(x)$, the second exposure column interchanges $1/2$ and $1/4$, and the action cost is
\[
 R(a)=a_1^2+a_2^2+4(a_1^4+a_2^4)+a_1a_2/4.
\]
An independent common-direction uniform innovation of radius $1/64$ is added. This is the previously specified extension, nested at $a_2=0$ when that innovation is disabled. Both models preserve $[0,1]^d$. Let $m\in\{1,2\}$ denote the number of controls and define
\begin{align*}
 Q_t^*(x,a)&=S(x)+R(a)+\beta\mathbb E V_{t+1}^*(F(x,a,Z)),\\
 V_t^*(x)&=\min_{a\in A(x)}Q_t^*(x,a),\qquad V_T^*=g.
\end{align*}
All costs and laws in this display are the economic primitives, not a fitted surrogate.

\begin{theorem}[Regularity of the original Bellman problem]\label{thm:regularity62}
For either control specification and every finite horizon, $V_t^*$ is convex and continuously differentiable in the interior, with the corresponding one-sided extension on the cube. It is $G_t$-Lipschitz in the $\ell^1$ norm and $M_t$-semiconcave in the Euclidean norm, where valid constants are
\begin{align*}
 G_T&=10/d,&G_t&=\frac{1037}{128d}+\beta\frac{47}{64}G_{t+1},\\
 M_T&=26/d,&M_t&=\frac{22}{d}+\frac{21}{256d}
             +\beta\left(\frac9{16}M_{t+1}+\frac18G_{t+1}\right).
\end{align*}
In particular, $G_t\le27/d$ and $M_t\le54/d$. The function $Q_t^*$ is jointly strongly convex on its convex feasible graph, with state curvature at least $107/(128d)$ and action curvature at least $7/4$.

A valid semiconcavity constant along each coordinate line is
\begin{align*}
 D_T&=\frac9d+\frac{16}{d^2},\\
 D_t&=\frac5d+\frac{16}{d^2}+\frac{21}{256d^2}
 +\beta M_{t+1}\left(\frac{85}{256}+\frac{23}{256d}\right).
\end{align*}
Consequently, for the positive multilinear interpolant $I_n$ on a mesh of spacing $1/n$,
\begin{equation}\label{eq:interpolation62}
 0\le I_nV_t^*(x)-V_t^*(x)\le\kappa_t,
 \qquad \kappa_t=\frac{dD_t}{8n^2}.
\end{equation}
\end{theorem}

\noindent The proof is in Supplement, Section~\ref{supp:proof-accuracy62-print-1}.


The theorem is a structural statement about these economic primitives. It is not a claim that every controlled diffusion, recursive preference specification, or game in the complete development edition has a convex Bellman problem. Nor does it require a signed-weight neural fit to preserve convexity. Convex approximation of dynamic programs has an established literature; the additional work here is to verify the needed curvature for the original nonlinear capacity-constrained model and connect it to the actual NBO actor and complete error account.

\subsection{Certified nodal Bellman bounds}
Suppose nodal arrays satisfy $l_{t+1}\le V_{t+1}^*\le u_{t+1}$ at all grid vertices. Equation~\eqref{eq:interpolation62} gives the functions $I_nl_{t+1}-\kappa_{t+1}$ and $I_nu_{t+1}$ as valid lower and upper continuations. The lower function need not itself be convex; convexity is used for the unknown true value and for the comparison argument, not asserted of arbitrary endpoint arrays.

Divide each uniform innovation component into $q$ equal bins and use its conditional mean. Convexity gives a lower quadrature bound for the true continuation. Semiconcavity gives an upper error
\begin{equation}\label{eq:quadrature62}
 \zeta_{t+1}=\frac{M_{t+1}}{6q^2}
                      \sum_{r=1}^{s}\rho_r^2\|s_r\|_2^2,
\end{equation}
where $s=1$ or $2$, $\rho_1=1/32$, and $\rho_2=1/64$ in the extension. The expectation remains the original continuous law. At the terminal date $g$ is evaluated directly, so no terminal interpolation allowance is needed.

At a vertex $v$, let $\mathcal A_A(v)$ be the capacity simplex lattice with spacing $c(v)/A$, including all its faces. For $m=1,2$, a minimizer can be represented by neighboring lattice vertices whose coordinate variances sum to at most $m c(v)^2/(4A^2)$. A valid action smoothness constant is
\[
 H_t=\frac{21}{4}+\beta M_{t+1}\|B\|_2^2,
 \quad \|B\|_2^2=\frac{5d}{32}\ (m=1),\quad
 \|B\|_2^2=\frac{9d}{32}\ (m=2).
\]
Taylor's inequality and that barycentric representation imply
\begin{equation}\label{eq:action-cover62}
 0\le\min_{a\in\mathcal A_A(v)}Q_t^*(v,a)-V_t^*(v)
       \le\delta_t:=\frac{mH_t}{128A^2}.
\end{equation}
The linear Taylor terms cancel even when the optimum lies on a face. This is an error against the continuous feasible optimum, not only against the finite action menu.

Let $q_t^-(v,a)$ and $q_t^+(v,a)$ include all interpolation, quadrature and arithmetic errors and enclose $Q_t^*(v,a)$. Then
\begin{equation}\label{eq:nodal62}
 l_t(v)=\max\{0,\min_{a\in\mathcal A_A(v)}q_t^-(v,a)-\delta_t\},
 \qquad
 u_t(v)=\min_{a\in\mathcal A_A(v)}q_t^+(v,a)
\end{equation}
are valid nodal Bellman bounds. The zero truncation uses the nonnegative economic costs. Extra neural or conventional proposals may improve the upper side but do not remove the lower action-cover allowance. Positive interpolation, averaging and minimum maps are monotone and nonexpansive; the separate allowances in \eqref{eq:interpolation62}--\eqref{eq:action-cover62} are therefore propagated, rather than being discarded because an endpoint calculation contains the truth.

\subsection{From a fitted witness to the complete implemented policy}
Let $a_t(v)\in A(v)$ be the actual action returned at each vertex by any generator. It may be selected by the exact neural witness, a conventional fitted continuation, or their common verified fallback. Write $w_v(x)$ for the multilinear weights and set $\bar\pi_t(x)=\sum_vw_v(x)a_t(v)$. Affineness of capacity makes this action feasible. Let $q_t^{\mathrm{sel},+}(v)$ be a verified upper bound on its original-optimum continuation value $Q_t^*(v,a_t(v))$, using the same lower and upper Bellman reference for all generators.

\begin{theorem}[Actual-policy, whole-domain certificate]\label{thm:policy62}
Set
\[
 e_t=\max_v\{q_t^{\mathrm{sel},+}(v)-l_t(v)\}+\kappa_t.
\]
Suppose the implemented policy $\widehat\pi_t$ is feasible and differs from $\bar\pi_t$ by at most $b_t$ in action $\ell^1$ norm, uniformly over the state cube. Define
\[
 \Lambda_t=\frac{13}{16}+\beta\frac{3d}{8}G_{t+1},\qquad
 D_T^\pi=0,\quad D_t^\pi=e_t+\Lambda_tb_t+\beta D_{t+1}^\pi.
\]
Then, at every true state and date,
\[
 0\le J_t^{\widehat\pi}(x)-V_t^*(x)\le D_t^\pi.
\]
This statement is conditional on the actual recorded policy, not on success of its fitting algorithm. It applies equally to a pure fitted policy and to a policy augmented by verified common candidates.
\end{theorem}

\noindent The proof is in Supplement, Section~\ref{supp:proof-accuracy62-print-2}.


In the execution, observation is rounded down at forty bits. An acquired state $y\le x$ has no greater capacity than the true state. Each interpolated action is then rounded down at twenty bits using the lower endpoint of its arithmetic enclosure. If $L_{ij}$ bounds the derivative of action coordinate $i$ with respect to state coordinate $j$, and $\eta$ bounds native interpolation error, a valid implementation allowance is
\[
 b_t\le 2^{-40}\sum_{i,j}L_{ij}
           +m(2^{-20}+2\eta+\eta_{\mathrm{out}}).
\]
Here $L_{ij}$ is computed from the maximum adjacent nodal difference times $n$, and $\eta_{\mathrm{out}}$ covers endpoint subtraction. All these numbers are recorded outward. This new barycentric acquired actor is explicitly distinguished from the earlier discontinuous cell-constant actor; an old actor is not silently replaced during a purported fixed-policy revalidation.

\subsection{Arithmetic, work, and scope}
The native interpolation kernel forms nonnegative tensor weights and a weighted sum of at most $2^d$ stored binary64 endpoints. It is compiled without reassociation or floating-point contraction. For $K=(3d+4)2^d+32$, an explicit conservative absolute allowance is
\[
 \eta=\frac{K2^{-53}}{1-K2^{-53}}\max_v|z_v|+K2^{-1022}.
\]
The final term conservatively covers underflow. Input coordinates, affine indexing, finite endpoints and the binary64 format are checked. Roundoff expansion and the ordinary interval cost arithmetic remain separate from the analytic approximation bounds. Exact rational fixtures compare the native sum with its mathematical tensor interpolant.

With $N=(n+1)^d$ vertices, $K_A=A+1$ scalar actions or $K_A=(A+1)(A+2)/2$ two-control actions, the reference calculation uses $O(TNK_Aq^s2^d d)$ elementary work under the explicit tensor implementation. Stored reference values use $O(TN)$ words and a policy uses $O(TNm)$ words. Streaming batches avoid a state-pair table. Proposal fitting, exact rational witness work, actor transfer, serialization, independent cost inference and compilation are additional costs, not absorbed into this bound. The analytic errors are second order in state, action and innovation resolution; the state-count and tensor-evaluation factors still depend exponentially on dimension. The result supplies a matched original-optimum comparison without asserting that these costs have disappeared.

\section{Frozen neural continuations and policy-cost precision}\label{sec:frozen72}
The optimal-advantage estimator in Section~\ref{sec:centered69} makes a policy comparison informative, but its conditional optimal-value queries are expensive. A fitted continuation can replace those queries without changing the estimand. The necessary certificate concerns its Bellman residual across states, not its distance from an unknown optimizer. This places the original NBO continuation, rather than only a proposed action, inside the evaluation calculation.

\subsection{Residual oscillation is the relevant range}
Consider the bounded expected-cost model with common terminal cost $g$, discounts $\beta_t\ge0$, and a common initial-state law. Write $b_0=1$ and $b_t=\prod_{j<t}\beta_j$. Fix bounded measurable continuations $h_0,\ldots,h_T$, conditionally on their construction. Define
\begin{align*}
 Q_t^h(x,a)&=c_t(x,a)+\beta_tP_t^a h_{t+1}(x),\\
 r_t^h(x)&=\inf_{a\in A_t(x)}Q_t^h(x,a)-h_t(x),\\
 a_t^{h,\pi}(x)&=Q_t^h(x,\pi_t(x))-\inf_{a\in A_t(x)}Q_t^h(x,a).
\end{align*}
The infimum is over the original continuous feasible set. Suppose verified constants give $0\le a_t^{h,\pi}\le\eta_t^\pi$, and likewise for $\nu$, throughout the true state domain. Let $\omega_t$ bound the oscillation of $r_t^h$, and $\omega_T$ that of $g-h_T$. Set
\[
 \Omega_h=\sum_{t<T}b_t\omega_t+b_T\omega_T,
 \qquad H_h^\pi=\sum_{t<T}b_t\eta_t^\pi.
\]
These constants include acquisition, action rounding and arithmetic allowances wherever those operations affect the implemented policy or the verified functions. A small empirical residual alone is not such a certificate.

\begin{theorem}[Frozen-continuation comparison]\label{thm:frozen72}
Along the actual trajectory of $\pi$, put
\begin{equation}\label{eq:frozen-observation72}
 Z_h^\pi=\sum_{t<T}b_t\bigl[Q_t^h(X_t^\pi,\pi_t(X_t^\pi))-h_t(X_t^\pi)\bigr]
                 +b_T(g-h_T)(X_T^\pi).
\end{equation}
For any coupling of the two original trajectory laws, $D_h=Z_h^\pi-Z_h^\nu$ satisfies
\begin{align}
 \mathbb E D_h&=\mathbb E[J_0^\pi-J_0^\nu],\label{eq:frozen-identity72}\\
 -\Omega_h-H_h^\nu&\le D_h\le\Omega_h+H_h^\pi.\label{eq:frozen-range72}
\end{align}
Moreover, at every initial state,
\begin{equation}\label{eq:frozen-regret72}
 0\le J_0^\pi-V_0^*\le\Omega_h+H_h^\pi.
\end{equation}
Adding an arbitrary state-independent constant to each $h_t$ changes neither $D_h$ nor these oscillation bounds. With $h_t=V_t^*$, the result reduces to optimal-advantage centering.
\end{theorem}
\noindent The proof is in Supplement, Section~\ref{supp:frozen72}.

The terminal correction in \eqref{eq:frozen-observation72} is essential unless $h_T=g$. Conditional expectation, not pathwise equality with the realized cash-cost difference, supplies \eqref{eq:frozen-identity72}. The result uses the actual kernels and costs; it assumes neither a correct statistical transition model estimated from data nor a convex neural fit. The pointwise policy certificate also applies conditionally after an adaptive choice among valid continuations. The independent-sample confidence statement below requires its additional sampling conditions.

For $M$ conditionally fixed comparisons, suppose independent trajectory rows yield outward enclosures $[L_i,U_i]$ of $D_h$, intersected with its valid support. Write $R_h=2\Omega_h+H_h^\pi+H_h^\nu$. A simultaneous interval with failure probability at most $\alpha$ is
\begin{equation}\label{eq:frozen-ci72}
 \left[\overline L-R_h\sqrt{\frac{\log(2M/\alpha)}{2n}},\quad
       \overline U+R_h\sqrt{\frac{\log(2M/\alpha)}{2n}}\right].
\end{equation}
This follows by the bounded-sum inequality applied to both endpoint samples and a union bound. The numerical and sampling width allocations of Corollary~\ref{cor:precision69} apply with $R_h$. Fitting, uniform residual qualification, conditional integration and any subsequent selection are distinct paid operations. In particular, a cache that adapts across rows changes the controller's history: its rows cannot be treated as independent evaluations of one frozen Markov policy. A fixed cache or an independent episode-reset design is sufficient to retain the stated sampling construction.

\subsection{Conditional integration of an actual neural continuation}
A shallow rational affine--ReLU continuation has the form
\[
 h(y)=c+\ell^\top y+\sum_{j=1}^K v_j(w_j^\top y+u_j)_+.
\]
Its readout weights may have either sign. In the original investment laws the next state, conditional on $(x,a)$, is affine in one or two independent centered uniform innovations. A neuron therefore has projected center $z$ and nonnegative projection radii $A,B$. Write $z_+=\max\{0,z\}$.

\begin{proposition}[Exact conditional ReLU moments]\label{prop:moment72}
For independent uniform projections of radii $A,B>0$, the conditional means of a neuron are
\begin{align}
 \mathcal M_1(z;A)&=\frac{(z+A)_+^2-(z-A)_+^2}{4A},\label{eq:moment-one72}\\
 \mathcal M_2(z;A,B)&=\frac{1}{24AB}\bigl[(z+A+B)_+^3-(z+A-B)_+^3\nonumber\\
 &\hspace{37mm}-(z-A+B)_+^3+(z-A-B)_+^3\bigr].\label{eq:moment-two72}
\end{align}
A zero projection radius is removed; if all radii vanish the mean is $z_+$. Consequently $P^ah$ is a finite exact rational calculation for rational primitives, states, actions and weights. For $n$ trajectories and at most $K$ neurons per date, its arithmetic-operation count is $O(nTKd)$, exclusive of fitting, trajectory construction and uniform qualification. No optimal-value descendant query occurs in this calculation.
\end{proposition}
\noindent The proof is in Supplement, Section~\ref{supp:moments72}.

The operation count does not bound rational bit complexity. Near-zero projected radii can also make a naive floating evaluation unstable; directed rounding or rational evaluation, including the zero-radius branches, is required. The executed moment catalogue checks the rational identities against an independently written polygon integral. It is a conditional-integration experiment, not a new trained-policy cost comparison. The formulas do not imply that the achieved residual oscillation of an arbitrary fitted network is small.

\subsection{Qualifying a fitted continuation without a Bellman tree}
Uniform qualification need not query $V^*$. Suppose $Q_t^h$ is Lipschitz in state and action with constants $L_{x,t},L_{a,t}$, and $h_t$ has constant $L_{h,t}$, all in $\ell^1$. In the investment capacity simplex, an action feasible at $y$ can be repaired continuously at $x$ with displacement at most $\kappa\|x-y\|_1$, where $\kappa=1/(8d)$. Thus a valid residual Lipschitz constant is
\[
 L_{r,t}=L_{x,t}+\kappa L_{a,t}+L_{h,t}.
\]
Let a finite state cover have fill radius $\rho$, and let $[l_j,u_j]$ enclose $r_t^h$ at every cover point. Then
\begin{equation}\label{eq:residual-cover72}
 \min_jl_j-L_{r,t}\rho\le r_t^h(x)
       \le\max_ju_j+L_{r,t}\rho
\end{equation}
throughout the cube. At a cover point, a feasible action mesh of fill radius $s$ gives a lower infimum bound by subtracting $L_{a,t}s$ from the smallest action-value lower endpoint; the smallest upper endpoint is an upper bound. Exact moments can supply those action endpoints. None of these steps requires convexity of $h$ or a true-value curvature minorant for $Q^h$.

For $N_t$ state-cover points and $A_t$ actions per point, this qualification costs $O(\sum_tN_tA_tK_{t+1}d)$ arithmetic operations, plus certification of the moduli and any actor-specific greedy-gap bound. Its cover can still be expensive in high dimension. The result replaces a recursive optimal-value oracle by an explicit approximation-and-qualification tradeoff; it does not make that qualification free. The residual-cover proof, the action-mesh argument and the terminal allowance are given in the supplement. Value-function control variates have an established variance-reduction literature \citep{munos2006variance}; the additional requirement here is an original-law, whole-domain certificate for the fitted NBO continuation and the actual acquired policy.

\section{Fitting a continuation for qualification}\label{sec:qualification73}
A continuation used to evaluate a policy should be fitted to the error that determines the evaluation cost. The residual-cover bound in \eqref{eq:residual-cover72} supplies such an objective. We now make its optimization constructive. The future continuation is held fixed at each backward step; neither an optimal-value oracle nor success of unrestricted hidden-layer training is assumed.

Fix a date and suppress its subscript. Let $x_1,\ldots,x_N$ be a state cover of radius $\rho$ in the $\ell^1$ metric, and suppose
\[
 y_j^-\le (\mathcal T h_+)(x_j)\le y_j^+.
\]
These are verified bounds on the original continuous-action infimum, not only on a proposal menu. Write $h_\theta=\sum_{k=1}^K\theta_k\phi_k$, where the fixed features have verified Lipschitz constants $L_k$ and $|\theta_k|\le B_k$. The constant feature is allowed with $L_k=0$. Let $C$ be a verified Lipschitz constant of $\mathcal T h_+$, including state-dependent feasibility repair. With $F_{jk}=\phi_k(x_j)$, define
\begin{equation}\label{eq:readout-objective73}
 \mathcal W(\theta)=\max_j(y_j^+-F_j\theta)
       -\min_j(y_j^--F_j\theta)
       +2\rho\left(C+\sum_k L_k|\theta_k|\right).
\end{equation}

\begin{theorem}[Qualification-optimal readout]\label{thm:readout73}
For every feasible $\theta$, $\operatorname{osc}(\mathcal T h_+-h_\theta)\le\mathcal W(\theta)$. Minimizing \eqref{eq:readout-objective73} is a finite linear program with $2K+2$ variables and $2N+2K$ inequality constraints, apart from variable bounds. For probability vectors $p,q$ on the nodes, set $\lambda_k=2\rho L_k$ and
\begin{equation}\label{eq:readout-dual73}
 \mathcal D(p,q)=p^\top y^+-q^\top y^-+2\rho C
   -\sum_k B_k\bigl(|[F^\top(q-p)]_k|-\lambda_k\bigr)_+.
\end{equation}
Then $\mathcal D(p,q)\le\inf_\theta\mathcal W(\theta)$, with equality for some $p,q$. Thus any feasible readout and exact probability vectors certify an achieved optimization gap of at most $\mathcal W(\theta)-\mathcal D(p,q)$.
\end{theorem}

The proof is in the Technical Supplement. It constructs the primal constraints and dual support function explicitly. Rationalizing a solver's proposed readout and probability vectors gives a checkable certificate; a convergence flag is not the certificate. Label intervals, features, moduli, coefficient bounds and the cover remain part of the numerical contract.

Applying the theorem backwards gives a fitted continuation at every date. Its achieved ranges $\mathcal W_t$, together with the terminal oscillation and the actual implemented greedy-gap bounds, can be substituted directly into Theorem~\ref{thm:frozen72}. This connects the attained fitting objective to both original-optimum regret and the range of the original-law cost observation. Each readout minimizes its stated one-step qualification bound conditional on the fixed future. It does not minimize the joint multi-date range, the true policy loss, or a hidden-feature objective. The next backward step recomputes its labels and moduli from the continuation actually returned.

The label construction in \eqref{eq:residual-cover72} uses exact conditional moments for a shallow ReLU future and a verified continuous-action cover. Its work is charged in addition to the linear program, coefficient rationalization, modulus calculation and dual check. A conventional basis with the same number of coefficients receives the same labels and the same objective; it uses its own valid feature moduli. This is a comparison of qualified approximations, not an assertion that different bases span the same functions or that a neural encoding has an intrinsic advantage.

\section{Representation calibration and optimization-stable policies}\label{sec:calibration81}
The attained continuation in Section~\ref{sec:frozen72} is an economic object, not merely an optional action proposal. Its qualification must charge the represented function rather than an avoidable featurewise overestimate. We make this comparison explicit and then pass from an attained fit to a class of optimization outputs with the same implemented policy. All statements in this section retain the original two-date investment law, fixed terminal payoff, acquired cells and feasible action menus.

\subsection{A common affine calibration}
Write a two-state additive continuation as
\[
 h_\theta(x)=\theta_0+\sum_{j=1}^4\theta_{1j}\psi_j(x_1)
                         +\sum_{j=1}^4\theta_{2j}\psi_j(x_2).
\]
For the ReLU basis, $\psi_j(s)=(s-(j-1)/4)_+$. For the Bernstein basis, $\psi_j(s)=\binom4j s^j(1-s)^{4-j}$; the omitted zeroth coefficient on each coordinate is zero. Both representations have nine coefficients and use the same coefficient box. Define
\begin{align}
 L_R(\theta)&=\max_{i=1,2}\max_{k=1,\ldots,4}
                      \left|\sum_{j=1}^k\theta_{ij}\right|,\label{eq:relu-modulus81}\\
 L_B(\theta)&=4\max_{i=1,2}\max_{j=1,\ldots,4}
                      |\theta_{ij}-\theta_{i,j-1}|,
                      \qquad \theta_{i0}=0.\label{eq:bernstein-modulus81}
\end{align}
These are state Lipschitz moduli for the $\ell^1$ distance. The first is exact for the represented piecewise-affine function; the second is a derivative-control bound and need not be exact for a general polynomial.

Suppose the fixed-label Bellman value at nodes $x_i$ belongs to $[y_i^-,y_i^+]$, every state is within $\ell^1$ distance $\rho$ of a node, and $G$ bounds the Bellman target's state modulus. With $\phi_i^\top\theta=h_\theta(x_i)$, put
\begin{equation}\label{eq:qualification81}
 F(\theta)=\max_i(y_i^+-\phi_i^\top\theta)
       -\min_i(y_i^--\phi_i^\top\theta)+2\rho\{G+L(\theta)\}.
\end{equation}
Here the labels, $G$, the cover and the feature map are fixed before this readout is optimized. A change in the fitted future defines a new fixed-label problem.

\begin{proposition}[Affine-exact qualification]\label{prop:calibration81}
The moduli in \eqref{eq:relu-modulus81}--\eqref{eq:bernstein-modulus81} give valid residual enclosures and the residual-oscillation upper bound \eqref{eq:qualification81}. The objective is convex and admits a finite linear-programming representation. For every common affine continuation $h(x)=a+u x_1+v x_2$, whenever its two coefficient representations are feasible, both moduli equal $\max\{|u|,|v|\}$. Thus, for common labels, node cover and target modulus, the two qualification objectives agree on this affine function.
\end{proposition}

\begin{proof}
A ReLU coordinate has the displayed cumulative slopes on its four affine intervals. The dual norm of $\ell^1$ is $\ell^\infty$, so their largest coordinate magnitude gives the state modulus. A Bernstein derivative is a convex combination of $4(\theta_{ij}-\theta_{i,j-1})$ with degree-three Bernstein weights. This proves its bound. Residuals change by at most $(G+L)\rho$ between a node and a covered state, proving the enclosure and its oscillation bound. Each term in \eqref{eq:qualification81} is a maximum of finitely many affine functions. For an affine function, the ReLU coefficients are $(u,0,0,0)$ and $(v,0,0,0)$; the Bernstein coefficients are $(u/4,u/2,3u/4,u)$ and the analogous vector for $v$. All relevant slopes or derivative controls are $u$ or $v$.
\end{proof}

This calibration does not identify the two non-affine function classes, equalize every numerical conditioning effect, or remove all conservatism from Bernstein controls. It isolates a particular qualification difference. For example, the old coefficientwise modulus charges $ux_1$ as $|u|$ in its ReLU encoding but as $10|u|$ in its Bernstein encoding. The calibrated modulus charges both as $|u|$. The factorial experiment below retains the old fits and evaluates the new fits separately; the calibration is not an ex post correction of their recorded costs.

\subsection{Exact certificates for action stability}
Fix the terminal-label qualification problem and an attained coefficient vector $\widehat\theta$, with $U=F(\widehat\theta)$. On each first-date observation cell $C$, let $a_C$ be its attained robust greedy action, evaluated at the declared cell representative. The terminal policy is fixed by the common terminal payoff. For each different action $b$ in the same robust menu, exact conditional integration gives an affine contrast
\[
 g_{Cb}(\theta)=Q^{h_\theta}(x_C,b)-Q^{h_\theta}(x_C,a_C)
                       =c_{Cb}+d_{Cb}^\top\theta.
\]
The switch problem minimizes $F(\theta)$ on the coefficient box subject to $g_{Cb}(\theta)\le0$. Let $D_{Cb}$ be any verified lower bound on its infimum. If the constraint is infeasible, its infimum is $+\infty$. Define $D_C=\min_{b\ne a_C}D_{Cb}$ and
\begin{equation}\label{eq:stable-cells81}
 \mathcal S_\epsilon=\{C:D_C>U+\epsilon\},\qquad \epsilon\ge0.
\end{equation}
Strict inequality excludes ties as well as strict reversals.

\begin{theorem}[Optimization-to-policy transfer]\label{thm:stability81}
Every feasible readout satisfying $F(\theta)\le U+\epsilon$ selects exactly $a_C$ on every cell in $\mathcal S_\epsilon$. In particular this holds for every $\epsilon$-optimal solution of the fixed-label qualification problem, because its optimum is at most $U$.

Let $q_\epsilon$ be the initial-law probability of the cells outside $\mathcal S_\epsilon$. Suppose the total cash cost of every such complete feasible controller belongs to $[0,H]$. With the terminal policy fixed as above,
\begin{equation}\label{eq:cost-stability81}
 \sup_{F(\theta)\le U+\epsilon}
       |J^{\pi_\theta}-J^{\pi_{\widehat\theta}}|
                       \le Hq_\epsilon.
\end{equation}
If all cells are certified, every readout in that sublevel gives the same complete policy and the same original-law cost. Under the uniform initial law on $[0,1]^2$ and an $n$-by-$n$ equal-cell cover, $q_\epsilon$ is exactly the fraction of uncertified cells; boundary assignments do not change that probability.
\end{theorem}

\begin{proof}
If $b$ weakly displaced $a_C$, the readout would be feasible for its switch problem and hence $F(\theta)\ge D_{Cb}\ge D_C$. This contradicts $F(\theta)\le U+\epsilon<D_C$. Couple the attained and alternative controllers with the same initial state and innovations. On a certified initial cell their first actions agree, their next states agree, and their fixed terminal-policy actions agree. Their cash costs are identical there. Elsewhere the absolute cash-cost difference is at most $H$, giving \eqref{eq:cost-stability81} by integration.
\end{proof}

The theorem does not replace failed switch certificates by successful ones: a nonpositive surplus is unresolved. Nor does a terminal-label certificate automatically propagate through newly fitted label problems at a longer horizon. The numerical result here concerns exactly the two-date construction for which only $h_1$ determines the first-date action and the terminal actor is common.

The switch bounds have an exact dual representation. Write $L(\theta)=\max_j|s_j^\top\theta|$, coefficient limits $|\theta_k|\le B_k$, and $\alpha=2\rho$. For probability vectors $p,q$, nonnegative multipliers $\xi_j^+,\xi_j^-$ with $\sum_j(\xi_j^++\xi_j^-)\le\alpha$, and $\nu\ge0$, set
\begin{align}
 z&=\sum_i(q_i-p_i)\phi_i
       +\sum_j(\xi_j^+-\xi_j^-)s_j+\nu d_{Cb},\nonumber\\
 D(p,q,\xi,\nu)&=p^\top y^+-q^\top y^-+2\rho G
                        +\nu c_{Cb}-\sum_kB_k|z_k|.\label{eq:switch-dual81}
\end{align}
Every such number is a valid $D_{Cb}$. The proof and independent exact replay are given in Appendix~\ref{app:dual81}. A floating optimizer supplies proposed multipliers only; their nonnegativity, normalization, total mass and lower-bound value are verified rationally.

\subsection{Simultaneous economic conclusions for optimizer classes}
Suppose the frozen attained-policy comparison supplies a simultaneous interval $[l_{jk},u_{jk}]$ for $J^{\pi_j}-J^{\pi_k}$. Apply Theorem~\ref{thm:stability81} separately to their fixed-label sublevels, with radii $v_j=H_jq_j$ and $v_k=H_kq_k$. Then, on the same event,
\begin{equation}\label{eq:uniform-comparison81}
 J^{\pi_\theta}-J^{\pi_\vartheta}
       \in[l_{jk}-v_j-v_k,\ u_{jk}+v_j+v_k]
\end{equation}
simultaneously for every pair of readouts in those sublevels. No union bound over an uncountable optimizer class is needed: the extension is deterministic on the original finite-policy event. When both covers are entirely stable the interval is unchanged. Sharp controller-choice and deployment-allocation results in Section~\ref{sec:allocation79} can therefore be applied directly to these qualified optimizer classes after the stated enlargement. This transfers attained-cost evidence without identifying numerical fitting success with either population training reliability or a new economic observation.

\section{Sequential qualification of policy costs}\label{sec:sequential73}
A fixed-sample interval does not by itself justify repeated inspection until a controller appears acceptable. Nor do cost observations from changing controllers identify the last controller's cost. We extend the frozen-continuation identity to predictable episodes and distinguish these targets. This uses established time-uniform concentration principles \citep{howard2020time}; confidence sequences have also been used for gated deployment in contextual bandits \citep{karampatziakis2021}. Here the observations are verified, continuous-law NBO cost contrasts, not importance-weighted bandit rewards.

Let $\mathcal F_{i-1}$ contain all information before episode $i$. Conditional on it, choose two complete controllers, a common initial law, a continuation $h_i$ fixed through the episode, and valid coefficients of Theorem~\ref{thm:frozen72}. The resulting cost contrast $D_i$ has conditional mean $\Delta_i$ equal to the original-law expected cost difference and predictable support $[a_i,b_i]$. Let $[L_i,U_i]\subseteq[a_i,b_i]$ enclose $D_i$ for the entire continuous innovation cell. Write $V_n=\sum_{i\le n}(b_i-a_i)^2$. For fixed positive numbers $\lambda_j,w_j$, $1\le j\le J$, with $\sum_jw_j\le1$, define
\begin{equation}\label{eq:anytime-boundary73}
 B_n=\min_{1\le j\le J}
       \left\{\frac{\log(2/(\alpha w_j))}{\lambda_j}
                    +\frac{\lambda_j V_n}{8}\right\}.
\end{equation}

\begin{theorem}[Predictable original-law cost certificate]\label{thm:anytime73}
Under the preceding conditional-law, support and enclosure conditions, with probability at least $1-\alpha$, simultaneously for every $n\ge1$,
\begin{equation}\label{eq:anytime-interval73}
 \frac{\sum_{i\le n}L_i-B_n}{n}
 \ \le\ \frac1n\sum_{i\le n}\Delta_i
 \ \le\ \frac{\sum_{i\le n}U_i+B_n}{n}.
\end{equation}
Each endpoint may also be intersected with the corresponding average predictable support endpoint. The assertion holds at every almost surely finite stopping time. If the complete controller pair and initial-law/reset contract stay fixed, then $\Delta_i=\Delta$, even when the qualified continuation changes predictably between episodes. Otherwise the target is the average conditional cost difference, not the last controller's cost.
\end{theorem}

The proof applies conditional Hoeffding bounds to each endpoint and Ville's inequality to a finite collection of nonnegative supermartingales. It requires no independent or identically distributed deployment contexts. It does require the stated conditional transition laws. Within-episode cache adaptation is allowed only when the controllers are treated as history-dependent algorithms and the greedy-gap qualification covers every memory state that can occur. Repeating the same program identifier while carrying a different cache state does not establish a fixed conditional target.

The implementation uses a finite dyadic collection of $\lambda_j$ and exact rational upper bounds for the logarithms. These give conservative, computable versions of \eqref{eq:anytime-boundary73}. It reserves the support and target identity before accepting an episode. A failed numerical evaluation consumes the full previously qualified support interval and its actual work; it is not deleted after inspecting its outcome. A failure of the support certificate itself stops qualification.

\begin{corollary}[Costed replacement after optional stopping]\label{cor:adoption73}
For a fixed target pair, let $[l_N,u_N]$ be the interval at a finite stopping time $N$. Let a declared deployment count be $m\ge1$, a one-time replacement and qualification charge be $F\ge0$, and the incremental per-deployment charge be $c$, all in the same units as the policy cost. On the simultaneous event of Theorem~\ref{thm:anytime73}, the net incremental cost lies in
\[
 [l_N+F/m+c,\ u_N+F/m+c].
\]
Consequently $u_N+F/m+c<0$ certifies strict net improvement. Inclusion of $[l_N,u_N]$ in a declared equivalence band certifies equivalence at that band, not strict superiority. New adaptively selected controller pairs may be qualified in separate epochs with conditional error allocations whose sum is at most the overall error budget.
\end{corollary}

A wall-clock reading is not a monetary charge. Measured setup, validation, rejected evaluations and successful retrievals are retained separately; prices are explicit decision inputs. The number of qualification episodes is not automatically the number of deployments over which a setup cost is amortized. These distinctions let the original investment decision use a valid stopping rule without silently changing either its estimand or its units.

\noindent The following attained-cost experiment retains the legacy coefficientwise qualification rule. Section~\ref{sec:study81} subsequently tests calibrated qualification on a separately frozen factorial catalogue.
\section{Original-law qualification of constructed policies}\label{sec:cost77}
A continuation can have a verified fitting objective without its implemented policy having an identified economic advantage. The distinction is resolved by evaluating the same complete policy that the construction returns. We now connect the qualification-optimal readouts of Section~\ref{sec:qualification73} to independent original-law cost observations. Neither the fitted future nor the acquired actor is replaced during evaluation. The comparison concerns the constructed NBO policies, rather than a different recursive controller supplied with the same neural label.

\subsection{Conditional integration and numerical observations}
In the two-state, two-date investment economy, write $a_t=\pi_t(X_t)$ and $\beta=15/16$. A cash observation and a conditionally integrated observation are
\begin{align}
 C^\pi&=c(X_0,a_0)+\beta c(X_1,a_1)+\beta^2g(X_2),\label{eq:cash77}\\
 Y^\pi&=c(X_0,a_0)+\beta c(X_1,a_1)
                    +\beta^2\bar g(X_1,a_1),\label{eq:conditional77}
\end{align}
where $\bar g(x,a)=P^ag(x)$. If $m(x,a)$ is the deterministic part of the transition, the original common, opposite-signed innovation gives the exact identity
\begin{equation}\label{eq:last-moment77}
 \bar g(x,a)=g(m(x,a))+\frac{5}{3072}.
\end{equation}
The shortage term depends on the sum of the coordinates and is unchanged by that innovation. The quadratic state and difference terms supply the variance correction. Consequently,
\[
 \E Y^\pi=\E C^\pi=\E J_0^\pi(X_0),\qquad
 \operatorname{Var}(Y^\pi-Y^\sigma)
       \le\operatorname{Var}(C^\pi-C^\sigma)
\]
when both policies use the same initial state and innovations. The variance inequality concerns the exact observations; it does not assert that every finite-sample numerical interval must be narrower.

The state cube is invariant. On that cube, $0\le a\le1/4$, a stage cost is at most $103/64$, and the terminal cost is at most $37/16$. Both observations therefore lie in
\begin{equation}\label{eq:support77}
 [0,H],\qquad H=\frac{103}{64}(1+\beta)+\frac{37}{16}\beta^2
                  =\frac{21097}{4096}.
\end{equation}
A paired difference has the signed support $[-H,H]$. Negative common-path gains are not clipped. No original-optimum value query is required to compute either observation for a frozen acquired actor.

An initial or innovation bin is a continuous cell, not a point mass. The evaluator encloses every true state in the bin, every acquired cell reached by that state enclosure, and every resulting feasible action. When an enclosure crosses an acquisition boundary, all intersected actions enter the continuation enclosure. A second implementation evaluates the same primitives using integer-directed arithmetic and expanded polynomials. The reported endpoints are the hull of the two valid enclosures, intersected with the proved support. Arithmetic agreement is a check on the implementations; validity follows from the enclosure rules and the primitive identities, not from agreement alone.

\subsection{Precision and the decision criterion}
The following statement applies to any fixed finite collection of qualified policies, not only to the two feature bases used below. All construction information is fixed before the independent evaluation rows are drawn. Dependence among policies on a common row is allowed.

\begin{theorem}[Original-law cost qualification]\label{thm:precision77}
For each of $K$ declared estimands, let $Z_i$ be its original-law observation on independent row $i$, and suppose a deterministic full-bin calculation returns
\[
 a\le L_i\le Z_i\le U_i\le b,\qquad i=1,\ldots,N.
\]
The inclusion may instead hold for the original conditional expectation within the sampled bin. Put $R=b-a$, let $\bar L,\bar U$ denote sample means, and let $s_L^2,s_U^2$ be the unbiased endpoint sample variances. For $N\ge2$ and $\ell>0$, define
\[
 r_L=\sqrt{\frac{2s_L^2\ell}{N}}+\frac{7R\ell}{3(N-1)},\qquad
 r_U=\sqrt{\frac{2s_U^2\ell}{N}}+\frac{7R\ell}{3(N-1)}.
\]
With probability at least $1-4K e^{-\ell}$, simultaneously for the declared collection,
\begin{equation}\label{eq:interval77}
 \E Z_i\in[\bar L-r_L,\bar U+r_U].
\end{equation}
Before rounding the final endpoints, its width is exactly
\begin{equation}\label{eq:width77}
 (\bar U-\bar L)
 +\sqrt{\frac{2s_L^2\ell}{N}}+\sqrt{\frac{2s_U^2\ell}{N}}
 +\frac{14R\ell}{3(N-1)}.
\end{equation}
Upward square roots and outward final rounding preserve coverage. The first term includes all retained numerical and acquisition uncertainty; the middle terms reflect observed endpoint dispersion; the last term is the finite-support penalty. A declared full-width target is met only when their computed sum meets it.
\end{theorem}

The proof in the Technical Supplement applies the empirical-Bernstein inequality of \citet{maurer2009} to the lower and upper endpoint variables. The statistical inequality is established work. The contribution here is its application to the original-law costs of the actual constructed actors, with executable whole-bin observations, independent arithmetic and an explicit precision account. Refining only a simulation bin cannot remove the finite-support or sampling terms. Conversely, an inaccurate uniform regret certificate need not prevent an accurate comparison of two actual policy costs.

Precision, equivalence and improvement are different conclusions. For a cost-difference interval $[l,u]$ and a predeclared margin $\delta$, equivalence requires $-\delta<l\le u<\delta$. Noninferiority of the first policy requires $u<\delta$; strict superiority requires $u<0$. An interval can meet its width target and fail equivalence because it identifies a material difference. These decisions are made from exact endpoints, not from rounded table entries.

\begin{corollary}[Priced selection and adoption]\label{cor:purchase77}
On the simultaneous event of Theorem~\ref{thm:precision77}, let $[L_j,U_j]$ qualify the cost of policy $j$, and let $w_j$ be its declared nonnegative resource charge. A minimizer $\widehat j$ of $U_j+w_j$ satisfies, for every policy $j$,
\[
 J_{\widehat j}+w_{\widehat j}
       \le J_j+w_j+(U_j-L_j).
\]
The inequality holds for every assignment of resource prices on the same event. If the gain from replacing an installed controller is in $[l,u]$, then $m$ deployments, monetary scale $v\ge0$, qualification charge $q\ge0$ and replacement fee $f\ge0$ give expected net benefit in
\[
 [mv l-q-f,\;mv u-q-f].
\]
In particular, when $l>0$ and $v>0$, every $m>(q+f)/(vl)$ is certified beneficial.
\end{corollary}
\begin{proof}
The first assertion follows from $J_{\widehat j}\le U_{\widehat j}$, the minimizing property, and $U_j\le J_j+U_j-L_j$. These deterministic inequalities hold for all nonnegative prices once all cost intervals cover. The second assertion is multiplication by $mv$ and subtraction of $q+f$.
\end{proof}

The resource charges in this result must be specified, not inferred from normalized losses. An already paid construction and a prospective purchase have different charges. When only a shared qualification service has been timed, charging that entire service to a proposed replacement is conservative; inventing a per-policy marginal qualification time is not.

\subsection{The constructed-controller experiment}
The fixed catalogue consists of the four previously constructed NBO actors and the admissible zero-investment incumbent. ReLU and degree-four Bernstein continuations each have nine coefficients and are constructed on eight- or sixteen-cell coordinate covers. At a fixed cover, the two bases use the same acquisition contract and primitive information. The actor arrays, actual fitted futures and exact optimization certificates are unchanged. The zero-investment policy supplies an economically interpretable installed reference; it is not substituted for the strong conventional query competitors in the preceding study.

Before execution, the protocol fixes $N=524{,}288$, forty-bit continuous bins, the sample stream, the full-width target $1/256$, and the separate two-sided equivalence margin $1/256$. All five marginal costs and ten paired differences are reported for both observations in \eqref{eq:cash77}--\eqref{eq:conditional77}. Thus $K=30$ and $\ell=10$ give a joint failure bound below $1/100$. Reusing a row for another controller or estimator does not increase the independent sample count. The declared independence model concerns the simulated rows; a fixed pseudorandom seed establishes reproducibility, not mathematical independence.

\begin{table}[htbp]\centering
\caption{Original-law costs of the constructed controllers}\label{tab:cost77}
{\small\setlength{\tabcolsep}{4pt}\begin{tabular}{lcc}
\toprule
Controller & Expected cost & Gain over zero incumbent \\
\midrule
Zero & [0.773561, 0.780120] & -- \\
R8 & [0.607716, 0.613107] & [0.165305, 0.167552] \\
B8 & [0.684034, 0.690159] & [0.089044, 0.090443] \\
R16 & [0.607022, 0.612384] & [0.165995, 0.168280] \\
B16 & [0.682910, 0.689022] & [0.090168, 0.091580] \\
\bottomrule\end{tabular}}
\par\smallskip{\footnotesize R and B denote the original nine-coefficient ReLU and Bernstein constructions; the number is the coordinate cover size. All policies are frozen before the new evaluation rows. The displayed observations integrate the last innovation analytically without changing expected cost. Endpoints are rounded outward. Paired gain intervals, not differences of marginal cost intervals, determine gain signs.}
\end{table}

\begin{table}[htbp]\centering
\caption{Paired feature and acquisition comparisons}\label{tab:pairs77}
{\small\setlength{\tabcolsep}{4pt}\begin{tabular}{lcrc}
\toprule
First minus second & Expected-cost difference & Full width & Equivalent \\
\midrule
R8 $-$ B8 & [-0.077568, -0.075801] & 0.001767 & No \\
R16 $-$ B16 & [-0.077159, -0.075368] & 0.001791 & No \\
R8 $-$ R16 & [0.000227, 0.001191] & 0.000963 & Yes \\
B8 $-$ B16 & [0.000661, 0.001599] & 0.000937 & Yes \\
\bottomrule\end{tabular}}
\par\smallskip{\footnotesize A negative interval identifies lower cost for the first controller. Equivalence is determined from exact endpoints at the predeclared two-sided margin 1/256. The separate full-width target is also 1/256. Matched-basis comparisons hold the acquisition contract fixed; cross-resolution comparisons do not.}
\end{table}

\paragraph{Outcomes of the fixed precision design.}

Across both estimators, 20 of the twenty paired intervals meet the declared full-width target. Their widths range from 0.000936 to 0.002285. Exactly 4 meet the separate equivalence criterion. Neither a failed width target nor a failed equivalence comparison is omitted.

The conditionally integrated comparisons identify a positive gain over the zero-investment incumbent for 4 of the four constructed policies. These are gains of complete implemented controllers, not reductions in a qualification objective.

At the two matched acquisition resolutions, the ReLU-minus-Bernstein comparisons identify 2 lower-cost ReLU outcomes, 0 higher-cost outcomes, and 0 unresolved signs. This finite comparison is not a universal representation ranking.


The construction and the evaluation now refer to the same complete controllers. The resulting cost intervals, rather than a reduction in the residual-fitting objective, determine the adoption regions in Corollary~\ref{cor:purchase77}. Comparisons between the feature bases identify only what their paired intervals support. The earlier unresolved and adverse query-method comparisons remain distinct evidence and are preserved in the Numerical and Development Appendix.

Both interval implementations evaluate every new trajectory and every acquired-action rectangle. Additional rational midpoint checks test predetermined rows. The phase ledger separates controller loading, lookup preparation, binary64 evaluation, integer evaluation, exact statistical arithmetic, output and the complete parent process. The independent replay recomputes the entire new enclosure catalogue without replacing its original clocks. This evaluates the two-date constructed-policy problem; it is not an independent rerun of the older, much larger vector Bellman tree.

\section{Qualification, representation, and attained economic cost}\label{sec:study81}
The preceding attained-controller experiment used a coefficientwise qualification bound. Its ReLU--Bernstein ordering is an observation about those specified constructions. We now examine whether that ordering persists after the two bases receive the same qualification charge on their common affine functions. This directly tests a source of representation-dependent conservatism while preserving the original nonlinear investment economy and the Neural Bellman Operator construction.

\subsection{A prospective factorial comparison}
The design fixes two bases, two qualification rules and two acquisition resolutions. The eight constructed controllers and the installed zero controller are evaluated on the same original-law paths. Each basis has nine coefficients with absolute bounds eight, the same node covers, thirty-three-action search menus and downward action quantum $1/4096$. The calibrated rule replaces the coefficientwise modulus by \eqref{eq:relu-modulus81} or \eqref{eq:bernstein-modulus81}, both in residual qualification and in the regularity allowance passed to the previous fitted date. Both own-future continuations are constructed backwards. An exact primal--dual gap of at most $1/65536$ qualifies each new fit. The coefficient proposal is quantized at forty bits; no hidden-feature optimizer or outcome-selected basis is introduced.

A source manifest was committed before construction, and the complete actor identities were committed before generating 524,288 new four-coordinate continuous-bin rows. The bin precision is forty bits. Two estimators, cash cost and exact last-innovation conditional expectation, evaluate the same expected costs under the original uniform initial law and continuous innovations. For each estimator the frozen family contains nine marginal costs, thirty-six paired differences and two basis-by-qualification interactions. All ninety-four estimands share a single failure allowance bounded by $376\exp(-11)<1/100$. Shared paths and the two estimators do not multiply the number of independent rows.

The whole-bin enclosures are independently computed in outward binary64 arithmetic and directed-integer arithmetic. Their hull, intersected with the proved support, is used for inference. An exact rational implementation additionally checks 512 midpoint trajectories per controller. All nine controllers, bins, endpoints, exact sample moments and disagreements are part of the retained protocol. This revision re-evaluates every bin and every inference quantity; it neither regenerates those observations as a second sample nor substitutes replay time for the original service clock.

\begin{table}[htbp]
\centering
\caption{Basis comparison before and after qualification calibration}\label{tab:basis81}
{\small\setlength{\tabcolsep}{3.5pt}
\begin{tabular}{lrcc}
\toprule
Qualification & $n$ & ReLU minus Bernstein & Lower cost \\
\midrule
Legacy & 8 & [-0.077633, -0.075733] & R \\
Legacy & 16 & [-0.077232, -0.075307] & R \\
Calibrated & 8 & [0.003571, 0.004688] & B \\
Calibrated & 16 & [0.000955, 0.002013] & B \\
\bottomrule\end{tabular}}
\par\smallskip{\footnotesize All entries are simultaneous paired original-law cost intervals from the same prospective R80 family, with exact last-innovation conditional integration. Positive values favor Bernstein. Equal coefficient counts and affine-exact modulus charging do not make the two nonlinear function classes identical. The old fits are evaluated, not overwritten.}
\end{table}

\begin{table}[htbp]
\centering
\caption{Actual controller gains from calibrated qualification}\label{tab:calibration-gains81}
{\small\setlength{\tabcolsep}{3.5pt}
\begin{tabular}{lrc}
\toprule
Basis & $n$ & Legacy minus calibrated cost \\
\midrule
Relu & 8 & [0.000457, 0.001653] \\
Relu & 16 & [0.005159, 0.006341] \\
Bernstein & 8 & [0.080964, 0.082772] \\
Bernstein & 16 & [0.082601, 0.084405] \\
\bottomrule\end{tabular}}
\par\smallskip{\footnotesize Each comparison retains its basis, coefficient count and bounds, robust action menus, observation resolution and original economic primitives. Both attained future fits are reconstructed under the stated qualification rule. The displayed gains are not timing estimates or claims about unrestricted neural fitting.}
\end{table}


The calibration improves the expected cost of each basis at each resolution. In particular, the fine ReLU controller gains between $0.005159$ and $0.006341$ relative to its legacy construction. The matched basis ordering nevertheless reverses: calibrated Bernstein has lower cost than calibrated ReLU at both resolutions. At $n=16$ the ReLU-minus-Bernstein interval is $[0.000955,0.002013]$, rounded outward. That strict ordering coexists with equivalence at the declared margin $1/256$; a small positive difference can satisfy both statements. The corresponding coarse-resolution interval does not establish equivalence at that margin.

\begin{table}[htbp]
\centering
\caption{Predeclared basis-by-qualification interactions}\label{tab:interactions81}
{\small\setlength{\tabcolsep}{3.5pt}
\begin{tabular}{rc}
\toprule
$n$ & Change in Bernstein-minus-ReLU cost \\
\midrule
8 & [-0.082254, -0.079371] \\
16 & [-0.079209, -0.076297] \\
\bottomrule\end{tabular}}
\par\smallskip{\footnotesize The interaction is (calibrated Bernstein minus calibrated ReLU) minus (legacy Bernstein minus legacy ReLU). These four-term contrasts were included in the source-frozen family before the new cost rows were generated. They are not obtained by treating the four marginal estimates as independent.}
\end{table}


The predeclared interactions isolate the change in the basis comparison. They do not establish that the bases span identical functions, that the Bernstein derivative-control bound is globally exact, or that every other numerical difference has been removed. The positive operational conclusion is more specific: charging the function's derivative structure changes the attained original-law policy cost, and a policy ranking cannot be attributed to a representation name without examining its qualification rule. The historical ReLU gain and the new reversed matched ranking remain separate valid observations.

\subsection{Precision of the economic comparison}
The target was a paired total width of $1/256$, fixed before seeing the new rows. Appendix~\ref{app:factorial81} reports every interval and its width components, including every marginal interval whether or not its width falls below that paired target. Variance reduction by conditional integration affects the variance component, not the bounded-support penalty or the acquisition contract. For example, in the fine calibrated basis comparison the conditional interval's width consists of approximately $0.001008605$ from the finite-support term, $0.000048152$ from endpoint sampling-variance allowances, and $2.69\times10^{-11}$ from the numerical enclosure. The support term is therefore the principal remaining contribution in that comparison.

The numerical component includes the implemented handling of continuous-bin width, acquisition ambiguity and arithmetic propagation. It is not represented as a pure floating-point error estimate. Likewise, a small numerical component does not justify dropping the finite-sample support term. These diagnostics explain the achieved precision without inferring significance from raw mean differences or treating the inherited R67 policy-cost rows as observations of these newly constructed actors.

\subsection{From a fitted readout to a qualified optimizer class}
For each new terminal-label qualification problem the action-switch catalogue covers every first-date observation cell and all thirty-two alternative robust actions. The 20,480 rational dual bounds are independently recomputed using the antiderivative and box-infimum calculations of Appendix~\ref{app:dual81}; no switch optimization is rerun to replace an unfavorable certificate.

\begin{table}[htbp]
\centering
\caption{First-date action stability of qualification sublevels}\label{tab:stability81}
{\small\setlength{\tabcolsep}{3.5pt}
\begin{tabular}{lrrrrc}
\toprule
Actor & Cells & $\epsilon=0$ & $1/65536$ & $1/4096$ & Uniform allowance \\
\midrule
CR8 & 64 & 28 & 28 & 28 & -- \\
CB8 & 64 & 64 & 64 & 55 & $1/32768$ \\
CR16 & 256 & 256 & 248 & 195 & $1/262144$ \\
CB16 & 256 & 256 & 243 & 200 & $1/8388608$ \\
\bottomrule\end{tabular}}
\par\smallskip{\footnotesize Counts use the strict inequality D greater than U plus epsilon. U is the exact objective of the attained feasible reference fit. The last column is a dyadic allowance strictly below the smallest positive switch surplus, when such a surplus is certified. A dash is unresolved, not evidence of instability. The independent verifier checks every one of the 20,480 action-switch duals.}
\end{table}

\begin{table}[htbp]
\centering
\caption{Cost comparisons uniform over exact qualification minimizers}\label{tab:optimizer-classes81}
{\small\setlength{\tabcolsep}{3.5pt}
\begin{tabular}{lc}
\toprule
Complete optimizer classes & Original-law cost difference \\
\midrule
CB16 minus CB8 & [-0.003286, -0.002241] \\
CB16 minus CR16 & [-0.002013, -0.000955] \\
CB8 minus CR16 & [0.000748, 0.001812] \\
\bottomrule\end{tabular}}
\par\smallskip{\footnotesize CB8, CR16 and CB16 have fully stable first-date covers at epsilon zero. Their terminal actors are fixed by the common terminal payoff. Consequently each stated exact-minimizer class gives the attained complete policy and inherits its interval without enlargement. This is a deterministic transfer of the same finite statistical family, not new sampling over optimizers.}
\end{table}


At zero extra objective allowance, the two fine-resolution constructions and coarse Bernstein have complete first-date action stability. Every exact minimizer of each of these fixed-label qualification problems therefore gives its attained complete policy, not merely a similar fitted value. The original-law intervals in Table~\ref{tab:optimizer-classes81} extend to these optimizer classes without additional probability expenditure. The same conclusion holds for the positive dyadic sublevels displayed in Table~\ref{tab:stability81}.

The standard $1/65536$ allowance leaves some fine-grid cells unresolved. Theorem~\ref{thm:stability81} then supplies an explicit enlargement based on their uniform initial mass. This conservative enlargement need not preserve a small cross-method sign. Coarse ReLU also retains unresolved cells even at zero allowance. None is relabeled as a discovered instability, and exact-minimizer invariance is not a claim that an arbitrary fitting routine reaches the requisite sublevel. The new statement addresses optimizer-independent policy stability under explicit attained-objective conditions; it does not estimate a hidden-layer success probability.

\subsection{Controller choice and complete costs}
Both cash and conditional marginal and pair intervals can be used together on their original simultaneous event. Exact shortest-path closure gives consistent cost potentials and the pure-choice regret bounds of the inherited decision theorem. The four-term interactions remain separately reported; sharpness of the pair-constraint calculation is not asserted after adding those extra restrictions. The complete nine-controller closure and attaining potentials are deposited in the revision audit.

For the fine calibrated controllers, the positive ReLU-minus-Bernstein interval determines a nonempty region of additional Bernstein charges for which its lower policy cost outweighs its extra charge. It also determines the opposite region in which ReLU is preferred after that charge. All such statements are conditional on an explicitly supplied conversion of computation, installation and deployment into normalized cost units. The exact boundary values may be tightened by the common pair-constraint closure; no new simulation is needed for that deterministic calculation.

The construction work is recorded rather than presumed equal. Legacy clocks include a label-loss ablation that the calibrated services do not execute. Their difference therefore does not identify an end-to-end speed gain from calibration. Exact action and expected-feature evaluation counts are reported alongside the clocks, and the original cost-inference, barrier and publication services retain separate measured costs. The present result supplies an identified policy-cost and conditional purchase comparison in the original economy. It does not extrapolate a two-date fitted-continuation calculation into a measured long-horizon speed result for the recursively queried vector tree.

\section{Deployment choice under qualified cost uncertainty}\label{sec:allocation79}
The attained continuation is an input to a complete economic decision, not merely to another approximation score. The same qualified costs can support a choice of acquisition precision across deployments. We retain the original NBO actors and draw a controller before observing the initial state or innovations. This implements a mixture of complete policies, not an average of their actions, coefficients or information sets. Its expected original-law cost is the corresponding mixture of the qualified costs. Minimax-regret allocation has an established decision-theoretic interpretation \citep{manski2004}; the result below concerns the finite cost set supplied by NBO qualification, conditional on its simultaneous coverage event.

Let $\mathcal C$ and $d(j,i)$ be the consistent difference-constraint set and distances of Theorem~\ref{thm:choice78}, stated and proved in the Numerical and Development Appendix. For a probability vector $p\in\Delta_q$ and per-deployment charges $w$, set $H_{ji}(w)=d(j,i)+w_i-w_j$, for $i,j=1,\ldots,q$. The charges include all costs that actually vary linearly with deployment shares; additional installation costs are treated separately.

\begin{theorem}[Sharp deployment allocation]\label{thm:allocation79}
The largest expected regret of allocating complete controllers with probabilities $p$ is
\begin{equation}\label{eq:allocation79}
 \mathcal R(p,w)=\max_{\theta\in\mathcal C}
 \left\{\sum_i p_i(\theta_i+w_i)-\min_j(\theta_j+w_j)\right\}
 =\max_j\sum_i p_iH_{ji}(w).
\end{equation}
Thus the optimal allocation and its value $V(w)$ are the solution of a finite linear program. For any rational $p,\lambda\in\Delta_q$, the computable numbers
\[
 L=\min_i\sum_j\lambda_jH_{ji}(w),\qquad
 U=\max_j\sum_i p_iH_{ji}(w)
\]
satisfy $L\le V(w)\le U$. There exist optimal primal and dual vectors with $L=U$. On the original coverage event, every implemented allocation has expected catalogue regret at most $U$, simultaneously for all $p,w$, including choices made after examining the qualified intervals. The worst-case bound in \eqref{eq:allocation79} is attained by one of the cost vectors already constructed in Theorem~\ref{thm:choice78}.
\end{theorem}

This is a decision calculation from the purchased observations, not a new experiment or an assertion that uncertainty reduction changes the true controller costs. Pure choices remain available as vertices of $\Delta_q$. Consequently $V(w)\le\min_i\Gamma_i(w)$, with strict improvement possible when no pure controller is uniformly best. The comparator is the best single qualified controller at the true cost vector. The original Bellman-optimum guarantees and the empirical comparisons between fitted representations remain separate.

\paragraph{An exact acquisition-fee region.}
Let $A$ and $B$ denote the fine and coarse ReLU controllers. The qualified gain is $J_B-J_A\in[l,u]$, where $0<l<u$. All other controllers are dominated by $B$ when their charges are equal. Suppose only $A$ has the incremental per-deployment charge $f$. For $l<f<u$, the fine-controller share and sharp expected regret are
\begin{equation}\label{eq:two-allocation79}
 p_A^*(f)=\frac{u-f}{u-l},\qquad
 V(f)=\frac{(f-l)(u-f)}{u-l}.
\end{equation}
For $f\le l$, use $A$ alone; for $f\ge u$, use $B$ alone. At the fee midpoint, allocation halves the best pure worst-case regret: $(u-l)/4$ rather than $(u-l)/2$. This is an exact decision bound, not a newly measured expected-cost reduction. In the retained experiment these bounds are approximately $0.0002405736$ and $0.0004811472$, respectively; Table~\ref{tab:allocation79} gives outward endpoints from the original rational intervals.

Maintaining both controllers need not be free. If a nontrivial allocation incurs an additional normalized cost $\kappa$, not incurred by either pure choice, its advantage survives precisely when
\begin{equation}\label{eq:setup-threshold79}
 \kappa<\frac{\min\{f-l,u-f\}^{2}}{u-l}.
\end{equation}
The full installed-subset account and a deterministic finite-cohort implementation are proved in Appendix~\ref{app:allocation79}. They charge every selected controller and do not turn shared experimental qualification into an unmeasured cold-start saving. For $m$ deployments and monetary loss scale $v>0$, the two-controller setup allowance is $mv$ times the right side. This gives an economically explicit frontier conditional on the actual fee, loss scale and deployment count, without inventing their calibration.

\section{Conclusion}\label{sec:conclusion65}
Neural Bellman Operators connect attained continuation construction to feasible policy decisions under the original law. The original witness and residual-qualification results provide the policy-loss account. Affine-calibrated derivative bounds make a representation-dependent qualification charge explicit, and rational action-switch barriers extend attained-policy conclusions to stated optimization sublevels.

The prospective factorial comparison identifies policy-cost gains from calibrated qualification in both bases. It also reverses the earlier matched basis ranking, demonstrating why equal feature counts alone do not identify a representation effect. Complete action-stable classes inherit their attained-cost intervals without additional sampling; unresolved cells retain their explicit enlargement. These constructive and economic conclusions do not require a universal neural ordering, a hidden-layer success assumption, or an unmeasured long-horizon speed claim.

The conditional purchase and deployment results use the qualified costs and declared resource charges as separate inputs. Original construction, failed reuse, adverse query timings, vector descendant growth and broader applications remain available with their own assumptions and records. A numerical decision is useful only through the complete chain from an attained approximation to a qualified and economically interpretable policy.

\bibliographystyle{ecta-fullname}
\bibliography{references,references78}
\end{document}
